\section{Relationships, cooperation, and conflict}\label{sec:peter-relations}

Peter's vocation was never an isolated ascent. He is almost always located in
a web of kinship, discipleship, collegial service, rivalry, correction, and
reconciliation. That web guards against two distortions: making him an
autocrat above the apostolic body, or making him merely an interchangeable
spokesman with no distinctive responsibility.

\subsection{Andrew, the Twelve, and the beloved disciple}

Andrew is Peter's brother and, in John, the one who brings him to Jesus. The
Synoptics more often place Peter with James and John. Their intimacy does not
eliminate competition: they dispute greatness, and James and John seek places
of honor. Peter asks what reward will come to those who have left everything.
Jesus repeatedly redirects rank toward service and suffering.

John's final chapters pair Peter with the beloved disciple. They run to the
tomb; Peter enters first, while the other disciple is said to see and believe.
At the lake, the beloved disciple recognizes Jesus and Peter jumps into the
water. After Peter's commission, he asks what will happen to the other, and
Jesus refuses the comparison: ``What concern is it of yours? You follow me.''
The narrative permits different forms of witness without forcing one to
cancel the other.

\subsection{James and the Jerusalem church}

Paul names Cephas, James, and John as reputed pillars who extended the right
hand of fellowship to Paul and Barnabas. The order differs elsewhere: at his
first Jerusalem visit Paul sees Cephas and James; Acts increasingly presents
James, the brother of the Lord, as the settled leader in Jerusalem. After his
escape from Herod, Peter tells the household to report to ``James and the
brothers'' and departs for another place. At the Acts 15 meeting, Peter gives
testimony, Paul and Barnabas report signs, and James voices the judgment.

No need exists to invent a struggle in which James displaced Peter, or a
constitutional chart in which every action flowed through Peter. The evidence
fits differentiated leadership: Peter bears a singular apostolic and
missionary prominence; James exercises a central Jerusalem role; the elders,
apostles, and whole church deliberate together.

\subsection{Paul: fellowship and public rebuke}

The relationship with Paul is among the best-attested and most contested.
Paul went to Jerusalem to visit or become acquainted with Cephas and stayed
fifteen days. Whatever else the verb in Galatians 1:18 implies, the meeting
joined the new apostle to the foremost Resurrection witness. At the later
meeting Cephas, James, and John recognized the grace given to Paul and agreed
upon differentiated missions: Paul and Barnabas toward the Gentiles, the
pillars toward the circumcised, with a shared charge to remember the poor.
This was communion, not an agreement never to cross an ethnic boundary. Peter
had Gentile contacts; Paul constantly addressed Jews and reasoned from
Israel's Scriptures.

Then Antioch exposed a fracture. Paul says Cephas ate with Gentiles but drew
back and separated himself after certain people came from James. Other Jewish
believers, including Barnabas, joined what Paul calls hypocrisy. Paul opposed
Cephas ``to his face'' before all and framed the issue through justification
and shared life in Christ. Galatians does not record Peter's response, a formal
settlement, or a permanent break.

The incident should neither be minimized into a harmless misunderstanding nor
inflated into proof of incompatible religions. Shared meals enacted whether
baptized Gentiles were full members without becoming Jews. Peter's withdrawal
therefore contradicted a truth his previous practice acknowledged. Catholic
tradition has never needed to deny the rebuke in order to confess Peter's
office. In fact, the episode distinguishes a charism of ecclesial service from
personal impeccability and shows apostolic correction occurring within
communion.

First Corinthians supplies later indirect evidence. Some Corinthians said
``I belong to Cephas,'' but Paul rejects every faction, including one bearing
his own name. His reference to Cephas's right to travel with a believing wife
assumes a recognized apostolate, not an enemy. The canonical letter called
Second Peter later refers to ``our beloved brother Paul'' and to his letters,
some passages of which the ignorant distort. Whatever judgment is made about
that letter's direct authorship, its canonical form receives Peter and Paul as
concordant witnesses without erasing the difficulty of Paul's writings.

\subsection{Mark and the problem of mediated memory}

First Peter 5:13 calls Mark ``my son.'' Papias, as quoted by Eusebius, says
that Mark became Peter's interpreter and accurately wrote remembered sayings
and deeds of the Lord, though not in order. Irenaeus develops a related
tradition. The notices make an early association between Peter, Mark, and the
Gospel tradition historically important. They do not turn the canonical
Gospel into a dictated autobiography or guarantee that every later claim
about Mark preserves Papias's meaning. Modern scholarship continues to debate
how much Petrine reminiscence can be identified behind Mark.

Peter's relationships thus reveal authority as received and answerable. He
strengthens brothers, accepts correction, testifies in councils, travels with
a wife, belongs among pillars, and is remembered through co-workers. The
``rock'' is never portrayed as a solitary Christian.
